
Or on a par
définition :
\[
 T_i=-\cl(\underline{J}/\underline{J}^2)=-N;\qquad
 T_j=-\cl(\underline{J}'/\underline{J}'^2)=-L.
\]
Par ailleurs, la suite exacte :
\[
0\longrightarrow \underline{F}\longrightarrow g^*(\underline{N})\longrightarrow \underline{O}_{Y'}(1)\longrightarrow0
\]
donne la relation :
\[
 g^*(T_i)=-(F+L),
\]
d'où on déduit :
\[
 \Todd(g^*(\check T_i))=\Todd(\check T_j)\cdot(\Todd(\check F))^{-1}.
\]

Compte tenu de 4.4, la relation à démontrer s'écrit donc :
\[
 \ch\bigl(j_*(g^*(x)\cdot\lambda_{-1}(F))\bigr)
 =
 j_{\mathrm{gr}}\bigl(\ch(g^*(x)\cdot\lambda_{-1}(F))\cdot\Todd(\check T_j)\bigr).
\]

Posant \(z=g^*(x)\cdot\lambda_{-1}(F)\), cette relation n'est autre que le théorème de Riemann--Roch appliqué à l'immersion régulière de codimension \(1\)~\(j\), et à l'élément \(z\) de \(K^{\bullet}(Y')\).

De plus, comme on a supposé que \(i\) vérifie l'hypothèse de (VII 4.1), \(\lambda_{-1}(F)\) est divisible par \((1-L)\). On peut donc trouver \(z'\) tel que \(z=z'(1-L)\) soit encore \(z=j^*(j_*(z'))\) (exp. VII 2.7). On est donc ramené au cas où \(Y\) est un diviseur, et à un élément \(x\) tel qu'il existe \(y\), avec \(x=i^*(y)\).

\textbf{4.6.} Il faut donc montrer la relation :
\[
 \ch\bigl(i_*(i^*(y))\bigr)
 =
 i_{\mathrm{gr}}\bigl(\ch(i^*(y))\cdot\Todd(\check T_i)\bigr).
\]
Or la formule de projection donne :
\[
 i_*(i^*(y))=y\cdot i_*(1),
\]
et de la suite exacte : \(0\longrightarrow \underline{J}\longrightarrow\underline{O}_X\longrightarrow\underline{O}_Y\longrightarrow0\), où \(\underline{J}\) est ici inversible, on tire :
\[
 i_*(1)=1-J,
\]
d'où
